\section{Патент на изобретение: Способ выявления многоэтапных кибератак в системе обнаружения вторжений Snort}

В рамках диссертационного исследования разработано техническое решение, направленное на повышение эффективности обнаружения сложных кибератак. Решение оформлено в виде заявки на патент на изобретение и программы для ЭВМ.

\subsection{Описание изобретения}

Изобретение относится к области информационной безопасности, в частности к системам обнаружения вторжений (IDS), и может быть использовано для выявления многоэтапных кибератак, включая целевые атаки (APT), на основе анализа последовательности событий с применением формальных моделей в виде атрибутивных метаграфов.

Способ заключается в перехвате событий срабатывания правил Snort, их агрегации во временном окне, построении атрибутивного метаграфа, сопоставлении с эталонными шаблонами MITRE ATT\&CK и генерации составного оповещения при превышении порога схожести.

\begin{figure}[h]
  \centering
  \includegraphics[width=0.95\textwidth]{patent_architecture.drawio}
  \caption{Архитектура системы обнаружения вторжений с модулем атрибутивных метаграфов}
  \label{fig:patent-arch}
\end{figure}

\begin{figure}[h]
  \centering
  \includegraphics[width=0.85\textwidth]{patent_metagraph_example.drawio}
  \caption{Пример атрибутивного метаграфа, построенного по событиям Snort}
  \label{fig:patent-metagraph}
\end{figure}

\subsection{Формула изобретения}

\begin{enumerate}[label=\arabic*., leftmargin=*]
  \item Способ выявления многоэтапных кибератак в системе обнаружения вторжений Snort, заключающийся в том, что:
  \begin{itemize}
    \item перехватывают события срабатывания правил Snort;
    \item агрегируют указанные события в заданном временном окне;
    \item формируют на основе агрегированных событий атрибутивный метаграф, в котором вершины соответствуют сетевым сущностям, а рёбра — связям между ними с атрибутами, включающими тактики и техники из базы знаний MITRE ATT\&CK, временные метки и уровень достоверности;
    \item сопоставляют сформированный метаграф с библиотекой эталонных метаграфов, описывающих известные цепочки компрометации;
    \item при превышении порогового значения меры схожести генерируют составное оповещение о многоэтапной кибератаке.
  \end{itemize}

  \item Способ по п.~1, отличающийся тем, что сопоставление осуществляют с использованием алгоритма частичного изоморфизма графов с учётом семантической близости атрибутов.

  \item Способ по п.~1, отличающийся тем, что составное оповещение передают в SIEM- или SOAR-систему через стандартный интерфейс вывода Snort.
\end{enumerate}

\subsection{Реферат}

Изобретение относится к области информационной безопасности и может быть использовано для обнаружения многоэтапных кибератак в системах обнаружения вторжений (IDS). Технический результат заключается в повышении полноты и точности выявления сложных, малозаметных атак за счёт интеграции в Snort модуля построения и анализа атрибутивных метаграфов. Способ включает перехват событий срабатывания правил Snort, агрегацию их во временном окне, формирование атрибутивного метаграфа, вершины которого соответствуют сетевым сущностям, а рёбра — связям с атрибутами MITRE ATT\&CK, сопоставление полученного метаграфа с эталонными шаблонами тактик компрометации и генерацию составного оповещения при превышении порога схожести. Изобретение обеспечивает обнаружение цепочек вида «фишинг → выполнение → боковое перемещение → утечка», не выявляемых стандартными сигнатурными механизмами, и может применяться в SOC, SIEM и системах защиты критической информационной инфраструктуры.

\begin{figure}[h]
  \centering
  \includegraphics[width=0.75\textwidth]{patent_workflow.drawio}
  \caption{Блок-схема алгоритма работы модуля MetaGraph-Snort}
  \label{fig:patent-workflow}
\end{figure}

\begin{figure}[h]
  \centering
  \includegraphics[width=0.9\textwidth]{patent_integration.drawio}
  \caption{Схема интеграции с SIEM/SOAR-платформами}
  \label{fig:patent-integration}
\end{figure}

\begin{figure}[h]
  \centering
  \includegraphics[width=0.8\textwidth]{patent_apt29_pattern.drawio}
  \caption{Эталонный шаблон атаки APT29 в виде атрибутивного метаграфа}
  \label{fig:patent-apt29}
\end{figure}